% Part: counterfactuals
% Chapter: introduction
% Section: material-conditional

\documentclass[../../../include/open-logic-section]{subfiles}

\begin{document}

\olfileid{cnt}{int}{mat}

\olsection{ભૌતિક શરતી વિધાન}

અંગ્રેજીમાં શરતી વાક્યનું સૌથી સરળ સ્વરૂપ ``જો \dots તો \dots''
હોય છે. અહીં \dots{}ની જગ્યાએ પોતે વાક્યો આવે છે; દાખલા તરીકે,
``જો બટલરે તે કર્યું હોય, તો માળી નિર્દોષ છે.'' પ્રારંભિક તર્કશાસ્ત્રના
અભ્યાસક્રમોમાં આપણે $\lif$ સંયોજક વડે શરતી વાક્યોનું પ્રતીકીકરણ
શીખીએ છીએ: \dots{}થી દર્શાવેલા ભાગોને, દાખલા તરીકે, !!{formula}s
$!A$ અને~$!B$ વડે પ્રતીકિત કરીએ છીએ; આખું શરતી વાક્ય
$!A \lif !B$ વડે દર્શાવીએ છીએ.

$\lif$ સંયોજક \emph{સત્યતાફલનલક્ષી} છે: $!A \lif !B$નું સત્યમૂલ્ય
---$\True$ કે $\False$---$!A$ અને~$!B$નાં સત્યમૂલ્યો પરથી નક્કી
થાય છે. $!A$ અસત્ય હોય અથવા $!B$ સત્ય હોય તો અને તો જ
$!A \lif !B$ સત્ય છે; અન્યથા તે અસત્ય છે. સત્યમૂલ્ય-નિયોજન
$\pAssign v$ના સંદર્ભમાં, $\pSat/{v}{!A}$ અથવા $\pSat{v}{!B}$ હોય
તો અને તો જ $\pSat{v}{!A \lif !B}$ એમ વ્યાખ્યાયિત કરીએ છીએ. આ
અર્થવિચાર ધરાવતા $\lif$ સંયોજકને \emph{ભૌતિક શરતી વિધાન} કહીએ છીએ.

આ વ્યાખ્યામાંથી કેટલાક પ્રાથમિક તાર્કિક તથ્યો મળે છે. સૌપ્રથમ,
ભૌતિક શરતી વિધાન માટે નિગમન પ્રમેય લાગુ પડે છે:
\begin{align}
  \text{જો } \Gamma, !A \Entails !B \text{ તો } \Gamma \Entails !A \lif !B
\end{align}
તે સત્યતાફલનલક્ષી છે: $!A \lif !B$ અને $\lnot !A \lor !B$ સમતુલ્ય છે:
\begin{align}
  !A \lif !B & \Entails \lnot !A \lor !B\\
  \lnot !A \lor !B & \Entails !A \lif !B
  \intertext{ભૌતિક શરતી વિધાન તેના પરિણામથી તથા તેના પૂર્વપક્ષના નિષેધથી ફલિત થાય છે:}
  !B & \Entails !A \lif !B\\
  \lnot !A & \Entails !A \lif !B
  \intertext{અસત્ય ભૌતિક શરતી વિધાન તેના પૂર્વપક્ષ અને પરિણામના નિષેધના સંયોજનને સમતુલ્ય છે:
    જો $!A \lif !B$ અસત્ય હોય, તો $!A \land \lnot !B$ સત્ય છે અને ઊલટું પણ ખરું:}
  \lnot(!A \lif !B) & \Entails !A \land \lnot !B\\
  !A \land \lnot !B & \Entails \lnot(!A \lif !B)
  \intertext{ભૌતિક શરતી વિધાન પર મોડસ પોનેન્સનો નિયમ લાગુ પડે છે:}
  !A, !A \lif !B & \Entails !B
  \intertext{ભૌતિક શરતી વિધાન પરિણામોને એકત્ર કરે છે:}
  !A \lif !B,  !A \lif !C & \Entails !A \lif (!B \land !C)
  \intertext{પૂર્વપક્ષને હંમેશાં વધુ મજબૂત કરી શકાય છે, એટલે કે આ શરતી વિધાન \emph{એકદિશવર્ધી} છે:}
  !A \lif !B & \Entails (!A \land !C) \lif !B
  \intertext{ભૌતિક શરતી વિધાન સંક્રામક છે, એટલે કે શૃંખલા-નિયમ માન્ય છે:}
  !A \lif !B, !B \lif !C & \Entails !A \lif !C
  \intertext{ભૌતિક શરતી વિધાન તેના પ્રતિપક્ષી સ્વરૂપ સાથે સમતુલ્ય છે:}
  !A \lif !B & \Entails \lnot !B \lif \lnot !A\\
  \lnot !B \lif \lnot !A & \Entails !A \lif !B
\end{align}

ગાણિતિક તર્કવિચારમાં આ બધાં ઉપયોગી અને નિર્વિવાદ અનુમાનો છે. પરંતુ
તત્ત્વચિંતન અને ભાષાવિજ્ઞાનના સાહિત્યમાં, બિનગાણિતિક સંદર્ભોમાં તેમને
અનુરૂપ અનુમાનો સામે મૂકાયેલા પ્રત્યુદાહરણો ઘણાં જોવા મળે છે. તે સૂચવે
છે કે અંગ્રેજીના ``જો \dots તો \dots'' વાક્યોનું પ્રતીકીકરણ કરતી વખતે
ભૌતિક શરતી વિધાન~$\lif$ યોગ્ય સંયોજક નથી---અથવા ઓછામાં ઓછું હંમેશાં
યોગ્ય નથી.

\end{document}
